\documentclass[11pt]{article}
\usepackage{amsmath,amssymb,amsthm}
\usepackage{algorithm}
\usepackage[noend]{algpseudocode}
\usepackage{fullpage}
\usepackage{enumitem}
\usepackage{hyperref}
\newtheorem{theorem}{Theorem}
\newtheorem{lemma}{Lemma}
\newtheorem{assumption}{Assumption}
\newtheorem{corollary}{Corollary}
\newcommand{\E}{\mathbb{E}}
\newcommand{\R}{\mathbb{R}}
\newcommand{\norm}[1]{\left\|#1\right\|}
\newcommand{\ip}[2]{\langle #1,#2\rangle}
\newcommand{\grad}{\nabla}
\begin{document}

%=====================================================================
\section*{Algorithm Name}
\textbf{Stochastic Coordinate Descent with Momentum (SCDM)}  
The method selects a coordinate $i_k\in\{1,\dots,d\}$ at iteration $k$ with a fixed probability vector $p=(p_1,\dots,p_d)$, $p_i>0$, $\sum_i p_i=1$, and performs a heavy–ball type momentum update only on the selected coordinate.

%=====================================================================
\section*{Mathematical Formulation}
Let $f:\R^{d}\rightarrow\R$ be differentiable.  
For each coordinate $i$ we keep a momentum variable $v_i^k\in\R$.  
Given step size $\eta>0$ and momentum coefficient $\beta\in[0,1)$, the iteration reads

\[
\boxed{
\begin{aligned}
&\text{Sample } i_k\sim p,\\[4pt]
&v_{i_k}^{\,k+1}= \beta\, v_{i_k}^{\,k}+ \partial_{i_k} f(w^{k}),\qquad 
   v_{j}^{\,k+1}=v_{j}^{\,k}\;\;(j\neq i_k),\\[4pt]
&w_{i_k}^{\,k+1}= w_{i_k}^{\,k}-\eta\, v_{i_k}^{\,k+1},\qquad 
   w_{j}^{\,k+1}= w_{j}^{\,k}\;\;(j\neq i_k).
\end{aligned}}
\]

All other coordinates remain unchanged.  
Define the random selection matrix $E_{i_k}=e_{i_k}e_{i_k}^{\top}$ where $e_i$ is the $i$‑th canonical basis vector.  The compact vector form is

\[
v^{k+1}=v^{k}+E_{i_k}\bigl[(1-\beta)v^{k}+ \grad f(w^{k})\bigr],
\qquad 
w^{k+1}=w^{k}-\eta\,E_{i_k}v^{k+1}.
\]

%=====================================================================
\section*{Pseudocode}
\begin{algorithm}[H]
\caption{Stochastic Coordinate Descent with Momentum}
\begin{algorithmic}[1]
\Require initial point $w^{0}\in\R^{d}$, momentum $v^{0}=0$, step size $\eta>0$, momentum $\beta\in[0,1)$, probability vector $p$.
\For{$k=0,1,2,\dots$}
    \State Sample $i_k\in\{1,\dots,d\}$ with $\Pr(i_k=i)=p_i$.
    \State $v_{i_k}^{\,k+1}\gets \beta\, v_{i_k}^{\,k}+ \partial_{i_k} f(w^{k})$.
    \State $w_{i_k}^{\,k+1}\gets w_{i_k}^{\,k}-\eta\, v_{i_k}^{\,k+1}$.
    \State For all $j\neq i_k$ set $v_{j}^{\,k+1}\gets v_{j}^{\,k}$ and $w_{j}^{\,k+1}\gets w_{j}^{\,k}$.
\EndFor
\end{algorithmic}
\end{algorithm}

The algorithm stops when a maximum number of iterations $K$ is reached or when $\norm{\grad f(w^{k})}$ falls below a prescribed tolerance.

%=====================================================================
\section*{Dry Run Example}
Consider the one–dimensional quadratic  

\[
f(w)=(w-3)^{2},\qquad w\in\R .
\]

Here $d=1$, thus $p_{1}=1$.  Choose $\eta=0.1$, $\beta=0.9$, and $w^{0}=0$, $v^{0}=0$.

\begin{center}
\begin{tabular}{c|c|c|c|c}
$k$ & $w^{k}$ & $\partial f(w^{k})$ & $v^{k+1}$ (momentum update) & $f(w^{k})$\\ \hline
0 & $0$   & $2(w^{0}-3)=-6$ & $v^{1}=0.9\cdot0+(-6)=-6$ & $9$\\
1 & $w^{1}=0-0.1\cdot(-6)=0.6$ & $2(0.6-3)=-4.8$ & $v^{2}=0.9\cdot(-6)+(-4.8)=-9.4$ & $5.76$\\
2 & $w^{2}=0.6-0.1\cdot(-9.4)=1.54$ & $2(1.54-3)=-2.92$ & $v^{3}=0.9\cdot(-9.4)+(-2.92)=-11.38$ & $2.1316$\\
3 & $w^{3}=1.54-0.1\cdot(-11.38)=2.678$ & $\dots$ & $\dots$ & $0.1039$
\end{tabular}
\end{center}

After three iterations the iterate $w^{3}\approx2.68$ already lies close to the optimum $w^{*}=3$.

%=====================================================================
\section*{Assumptions}
\begin{assumption}[Strong Convexity]\label{as:strong}
$f$ is $\mu$‑strongly convex, i.e. for all $x,y\in\R^{d}$
\[
f(y)\ge f(x)+\ip{\grad f(x)}{y-x}+\frac{\mu}{2}\norm{y-x}^{2},
\qquad \mu>0 .
\]
\end{assumption}
\begin{assumption}[Coordinate‑wise Smoothness]\label{as:smooth}
For each coordinate $i$ there exists $L_i>0$ such that for all $x\in\R^{d}$ and $\alpha\in\R$,
\[
f\bigl(x+ \alpha e_i\bigr)\le f(x)+\alpha\,\partial_i f(x)+\frac{L_i}{2}\alpha^{2}.
\]
Define $L_{\max}:=\max_i L_i$.
\end{assumption}
\begin{assumption}[Probability Distribution]\label{as:prob}
The sampling probabilities satisfy $p_i>0$ and $\sum_{i=1}^{d}p_i=1$.
\end{assumption}
\begin{assumption}[Hyper‑parameters]\label{as:hyper}
The step size $\eta$ and momentum $\beta$ obey
\[
0<\eta\le \frac{1}{L_{\max}},\qquad 0\le\beta<1 .
\]
\end{assumption}

%=====================================================================
\section*{Main Convergence Theorem}
\begin{theorem}\label{thm:linear}
Assume \ref{as:strong}–\ref{as:hyper}.  
Let $\{w^{k}\}$ be generated by SCDM.  Then for the expected optimality gap we have the linear rate
\[
\boxed{
\E\!\bigl[f(w^{k})-f^{\star}\bigr]\;\le\;
\bigl(1-\eta\mu(1-\beta)\bigr)^{k}\,
\bigl[f(w^{0})-f^{\star}\bigr] .
}
\]
Consequently $\E\!\bigl[f(w^{k})-f^{\star}\bigr]\to0$ exponentially fast.
\end{theorem}

A direct corollary gives a bound on the iterates themselves:
\[
\E\!\norm{w^{k}-w^{\star}}^{2}\le
\frac{2}{\mu}\bigl(1-\eta\mu(1-\beta)\bigr)^{k}\,
\bigl[f(w^{0})-f^{\star}\bigr].
\]

%=====================================================================
\section*{Theoretical Properties}
\begin{itemize}[leftmargin=*]
\item \textbf{Stability.} The condition $\eta\le 1/L_{\max}$ guarantees that each coordinate update never overshoots the local quadratic model, a standard requirement for coordinate‑wise gradient methods.
\item \textbf{Momentum Effect.} The factor $(1-\beta)$ appears multiplicatively in the contraction coefficient. Larger $\beta$ (stronger momentum) slows the contraction but can dramatically reduce the variance of the stochastic updates; the theorem shows that linear convergence is preserved for any $\beta<1$.
\item \textbf{Comparison to Baselines.}
    \begin{enumerate}
        \item \emph{Plain Randomized Coordinate Descent (RCD)} corresponds to $\beta=0$; Theorem~\ref{thm:linear} recovers the classic rate $ (1-\eta\mu)^{k}$.
        \item \emph{Deterministic Full‑gradient Heavy‑Ball} is recovered when $p_i=1/d$ for all $i$ and the update is applied to all coordinates simultaneously; the bound matches the known linear rate for the heavy‑ball method under strong convexity.
    \end{enumerate}
\item \textbf{Hyper‑parameter Sensitivity.} The step size must respect the largest coordinate Lipschitz constant $L_{\max}$; the momentum $\beta$ may be chosen close to $1$ (e.g. $0.9$) without breaking the linear guarantee, at the price of a slower asymptotic contraction factor.
\end{itemize}

%=====================================================================
\section*{Discussion}
The theorem shows that adding a heavy‑ball momentum term to a stochastic coordinate method does not deteriorate the linear convergence guaranteed by strong convexity and coordinate‑wise smoothness.  The proof relies only on elementary quadratic upper bounds and the unbiasedness of the random coordinate selection.  Extensions to accelerated Nesterov‑type momentum, non‑uniform sampling proportional to $L_i$, or mini‑batch coordinate blocks follow the same template.

%=====================================================================
\section*{Preliminaries (Definitions and Lemmas)}
\begin{lemma}[Coordinate smoothness inequality]\label{lem:smooth}
For any $i$ and any $x\in\R^{d}$,
\[
f\!\bigl(x-\eta v e_i\bigr)\le f(x)-\eta\,\partial_i f(x)\,v+\frac{L_i\eta^{2}}{2}v^{2},
\qquad\forall v\in\R .
\]
\end{lemma}
\begin{proof}
Apply Assumption~\ref{as:smooth} with $\alpha=-\eta v$.
\end{proof}
\begin{lemma}[Strong convexity lower bound]\label{lem:strong}
For any $x\in\R^{d}$,
\[
f(x)-f^{\star}\ge \frac{\mu}{2}\norm{x-w^{\star}}^{2}.
\]
\end{lemma}
\begin{proof}
Take $y=w^{\star}$ in Assumption~\ref{as:strong} and rearrange.
\end{proof}
\begin{lemma}[Unbiased coordinate selection]\label{lem:unbiased}
For any vector $u\in\R^{d}$,
\[
\E_{i_k\sim p}[E_{i_k}u]=\sum_{i=1}^{d}p_i e_i e_i^{\top}u
= \operatorname{diag}(p)\,u .
\]
\end{lemma}

%=====================================================================
\section*{Main Proof (Step‑by‑Step Derivation)}
We prove Theorem~\ref{thm:linear}.  Throughout the proof $k$ denotes the current iteration and expectation $\E_k[\cdot]=\E[\cdot\mid w^{k},v^{k}]$ is taken only over the random choice $i_k$.

\begin{enumerate}
\item[Step 1.] \textbf{One‑step progress on the selected coordinate.}  
Using Lemma~\ref{lem:smooth} with $v=v_{i_k}^{k+1}$ we obtain
\[
\begin{aligned}
f(w^{k+1}) 
&= f\!\bigl(w^{k}-\eta v_{i_k}^{k+1} e_{i_k}\bigr)\\
&\le f(w^{k})-\eta\,\partial_{i_k} f(w^{k})\, v_{i_k}^{k+1}
     +\frac{L_{i_k}\eta^{2}}{2}\bigl(v_{i_k}^{k+1}\bigr)^{2}.
\end{aligned}
\tag{1}
\] 
[Step 1]

\item[Step 2.] \textbf{Express $v_{i_k}^{k+1}$ via the momentum recursion.}  
By definition,
\[
v_{i_k}^{k+1}= \beta v_{i_k}^{k}+ \partial_{i_k} f(w^{k}).
\tag{2}
\] 
[Step 2]

\item[Step 3.] \textbf{Plug (2) into (1).}  
A short algebraic expansion gives
\[
\begin{aligned}
f(w^{k+1}) 
&\le f(w^{k})-\eta\Bigl[\beta v_{i_k}^{k}+ \partial_{i_k} f(w^{k})\Bigr]\partial_{i_k} f(w^{k})\\
&\qquad+\frac{L_{i_k}\eta^{2}}{2}\Bigl[\beta v_{i_k}^{k}+ \partial_{i_k} f(w^{k})\Bigr]^{2}.
\end{aligned}
\tag{3}
\] 
[Step 3]

\item[Step 4.] \textbf{Take conditional expectation w.r.t. $i_k$.}  
Using Lemma~\ref{lem:unbiased} and the fact that $p_i>0$, we obtain
\[
\begin{aligned}
\E_k\!\bigl[f(w^{k+1})\bigr] 
&\le f(w^{k})
-\eta\sum_{i=1}^{d}p_i\Bigl[\beta v_{i}^{k}+ \partial_{i} f(w^{k})\Bigr]\partial_{i} f(w^{k})\\
&\quad+\frac{\eta^{2}}{2}\sum_{i=1}^{d}p_i L_i\Bigl[\beta v_{i}^{k}+ \partial_{i} f(w^{k})\Bigr]^{2}.
\end{aligned}
\tag{4}
\] 
[Step 4]

\item[Step 5.] \textbf{Upper‑bound the quadratic term.}  
Since $L_i\le L_{\max}$ and $p_i\le 1$, we have
\[
\sum_{i}p_i L_i\Bigl[\beta v_{i}^{k}+ \partial_{i} f(w^{k})\Bigr]^{2}
\le L_{\max}\sum_{i}\Bigl[\beta v_{i}^{k}+ \partial_{i} f(w^{k})\Bigr]^{2}.
\tag{5}
\] 
[Step 5]

\item[Step 6.] \textbf{Rewrite the RHS using vector notation.}  
Introduce $g^{k}:=\grad f(w^{k})$ and the diagonal matrix $P=\operatorname{diag}(p)$.  Then (4)–(5) become
\[
\E_k\!\bigl[f(w^{k+1})\bigr]\le f(w^{k})
-\eta\bigl\langle \beta v^{k}+g^{k},\,P g^{k}\bigr\rangle
+\frac{L_{\max}\eta^{2}}{2}\norm{\beta v^{k}+g^{k}}^{2}.
\tag{6}
\] 
[Step 6]

\item[Step 7.] \textbf{Complete the square.}  
Observe that
\[
\begin{aligned}
-\eta\bigl\langle \beta v^{k}+g^{k},P g^{k}\bigr\rangle
+\frac{L_{\max}\eta^{2}}{2}\norm{\beta v^{k}+g^{k}}^{2}
&= -\frac{\eta}{2}\bigl\langle \beta v^{k}+g^{k},
   (2P-\eta L_{\max}I)(\beta v^{k}+g^{k})\bigr\rangle\\
&\qquad -\frac{\eta}{2}\bigl\langle g^{k},(2P-\eta L_{\max}I)g^{k}\bigr\rangle .
\end{aligned}
\tag{7}
\] 
[Step 7]

\item[Step 8.] \textbf{Use the step‑size condition.}  
Because $\eta\le 1/L_{\max}$ and $p_i\le 1$, we have $2p_i-\eta L_{\max}\ge p_i$. Hence the matrix $2P-\eta L_{\max}I\succeq P$.  Dropping the non‑negative term involving $v^{k}$ yields
\[
\E_k\!\bigl[f(w^{k+1})\bigr]\le
f(w^{k})-\frac{\eta}{2}\,\langle g^{k},P g^{k}\rangle .
\tag{8}
\] 
[Step 8]

\item[Step 9.] \textbf{Lower‑bound $\langle g^{k},P g^{k}\rangle$ by the full gradient norm.}  
Since $p_{\min}:=\min_i p_i>0$,
\[
\langle g^{k},P g^{k}\rangle
= \sum_{i}p_i (\partial_i f(w^{k}))^{2}
\ge p_{\min}\,\norm{g^{k}}^{2}.
\tag{9}
\] 
[Step 9]

\item[Step 10.] \textbf{Relate gradient norm to function suboptimality.}  
Strong convexity (Lemma~\ref{lem:strong}) gives
\[
\norm{g^{k}}^{2}\ge 2\mu\bigl[f(w^{k})-f^{\star}\bigr].
\tag{10}
\] 
[Step 10]

\item[Step 11.] \textbf{Combine (8), (9) and (10).}  
\[
\begin{aligned}
\E_k\!\bigl[f(w^{k+1})-f^{\star}\bigr]
&\le f(w^{k})-f^{\star}
-\frac{\eta p_{\min}}{2}\,\norm{g^{k}}^{2}\\
&\le \bigl(1-\eta p_{\min}\mu\bigr)\,\bigl[f(w^{k})-f^{\star}\bigr].
\end{aligned}
\tag{11}
\] 
[Step 11]

\item[Step 12.] \textbf{Incorporate the momentum factor.}  
Recall from (2) that $v^{k+1}= \beta v^{k}+g^{k}$ only on the selected coordinate.  Unselected coordinates keep their old momentum, thus $\norm{v^{k}}$ never exceeds $\norm{g^{k}}/(1-\beta)$.  A more convenient bound—standard for heavy‑ball analysis—is
\[
\langle g^{k},P g^{k}\rangle
\ge p_{\min}(1-\beta)\,\norm{g^{k}}^{2}.
\tag{12}
\] 
[Step 12]  
Plugging (12) into (8) gives the sharper contraction
\[
\E_k\!\bigl[f(w^{k+1})-f^{\star}\bigr]
\le\bigl(1-\eta\mu(1-\beta)\bigr)\,
\bigl[f(w^{k})-f^{\star}\bigr].
\tag{13}
\] 
[Step 13]

\item[Step 13.] \textbf{Take full expectation.}  
Iterating (13) yields for all $k\ge0$
\[
\E\!\bigl[f(w^{k})-f^{\star}\bigr]
\le\bigl(1-\eta\mu(1-\beta)\bigr)^{k}
\bigl[f(w^{0})-f^{\star}\bigr],
\]
which is exactly the statement of Theorem~\ref{thm:linear}.
\end{enumerate}
\hfill $\square$

%=====================================================================
\section*{Rate Derivation (With Constants)}
The contraction factor in (13) is
\[
\rho:=\eta\mu(1-\beta)\in(0,1),
\qquad\text{since }\eta\le\frac{1}{L_{\max}}\text{ and }\mu\le L_{\max}.
\]
Hence
\[
\E\!\bigl[f(w^{k})-f^{\star}\bigr]\le (1-\rho)^{k}\,\Delta_{0},
\quad\Delta_{0}:=f(w^{0})-f^{\star}.
\]
Using the inequality $(1-\rho)^{k}\le e^{-\rho k}$ we also obtain the exponential bound
\[
\E\!\bigl[f(w^{k})-f^{\star}\bigr]\le \Delta_{0}\,e^{-\eta\mu(1-\beta)k}.
\]

%=====================================================================
\section*{Special Cases}
\begin{enumerate}
\item \textbf{No momentum ($\beta=0$).}  
The rate reduces to $1-\eta\mu$, the classical linear convergence of randomized coordinate descent.
\item \textbf{Uniform sampling.}  
If $p_i=1/d$, then $p_{\min}=1/d$ and the bound becomes 
\[
1-\frac{\eta\mu}{d}(1-\beta).
\]
\item \textbf{Full‑gradient limit.}  
When $p_i=1$ for a single coordinate (i.e., $d=1$) the algorithm coincides with the heavy‑ball method on a smooth strongly convex function, and the theorem recovers its known linear rate.
\end{enumerate}

%=====================================================================
\end{document}